\documentclass[11pt,letterpaper]{article}
\usepackage[T1]{fontenc}
\usepackage[utf8]{inputenc}
\usepackage{lmodern}
\usepackage[spanish,mexico]{babel}
\usepackage[left=2cm,right=2cm,top=2.2cm,bottom=2.2cm]{geometry}
\usepackage[table]{xcolor}
\usepackage{array,booktabs,longtable,tabularx}
\usepackage{amsmath,amssymb}
\usepackage{fancyhdr}
\usepackage[hidelinks]{hyperref}

\definecolor{cR}{RGB}{31,94,168}
\definecolor{cM}{RGB}{46,125,50}
\definecolor{cO}{RGB}{123,31,162}
\definecolor{cC}{RGB}{176,84,0}
\definecolor{caja}{RGB}{243,246,250}
\definecolor{regla}{RGB}{200,200,200}

\newcommand{\etiq}[2]{{\setlength{\fboxsep}{1.4pt}\colorbox{#1}{\textcolor{white}{\sffamily\bfseries\scriptsize #2}}}}
\newcommand{\tR}{\etiq{cR}{RESULTADOS}}
\newcommand{\tM}{\etiq{cM}{MÉTODO}}
\newcommand{\tO}{\etiq{cO}{COMITÉ}}
\newcommand{\tC}{\etiq{cC}{CITA}}
\newcommand{\sep}{\\\arrayrulecolor{regla}\midrule\arrayrulecolor{black}}
\newcommand{\ra}{$\to$}
\newcolumntype{L}[1]{>{\raggedright\arraybackslash}p{#1}}

\setlength{\parindent}{0pt}
\setlength{\parskip}{5pt}
\setlength{\LTpre}{4pt}
\setlength{\LTpost}{6pt}
\renewcommand{\arraystretch}{1.15}

\pagestyle{fancy}
\fancyhf{}
\renewcommand{\headrulewidth}{0pt}
\fancyfoot[L]{\footnotesize\sffamily Cambios respecto a la versión revisada (9 de junio de 2026)}
\fancyfoot[R]{\footnotesize\sffamily \thepage}

\makeatletter
\renewcommand\section{\@startsection{section}{1}{\z@}{-14pt plus -2pt}{6pt}{\sffamily\large\bfseries}}
\renewcommand\subsection{\@startsection{subsection}{2}{\z@}{-10pt plus -2pt}{4pt}{\sffamily\normalsize\bfseries}}
\makeatother

\begin{document}

{\sffamily
{\LARGE\bfseries Cambios en la tesis respecto a la versión revisada}\par\vspace{4pt}
{\large Impresión revisada por el comité (9 de junio de 2026) frente a la versión actual (4 de noviembre de 2026)}\par\vspace{8pt}
\small
\textit{Logística de última milla en redes viales reales: un enfoque reproducible para el MDRP en La Paz, B.C.S.}\\
Alan José Marcos Sánchez Martínez \quad·\quad Director: Dr.\ Marco Antonio Castro Liera\\
Maestría en Sistemas Computacionales, Instituto Tecnológico de La Paz\par}

\vspace{6pt}
\noindent\colorbox{caja}{\parbox{\dimexpr\textwidth-2\fboxsep}{\vspace{4pt}
\textbf{\sffamily En resumen}
\begin{itemize}\setlength{\itemsep}{2pt}\setlength{\topsep}{2pt}
\item Se encontraron y corrigieron errores en el simulador (Sección~2). de lo mas destacable: el instante de entrega
  no incluía la espera de la comida en el restaurante ni los tiempos de servicio. Por eso se repeti la campaña
  experimental completa.
\item Con el simulador ya corregido, las diez semillas son válidas (antes, siete): la falla de tres semillas se debía
  a un error del simulador, no a una interrupción de OSRM. ahora hay 80 corridas en lugar de 56.
\item Los tiempos de entrega son mayores en ambas políticas y el punto de cruce se adelanta: RH supera a FCFS
  desde el nivel de saturación S2 (antes, desde S3). En S1, FCFS entrega en promedio 1.4~min antes, pero la diferencia en el cumplimiento del SLA
  (\textit{Service Level Agreement}; porcentaje de pedidos entregados en 40~min o menos) no es significativa.
\item Dado que no se entendia bien la idea principal del trabajo de inicio, cambie el planteamiento del
  problema (§1.2). reescribi el planteamiento con una pregunta central y tres preguntas de investigación, que se responden en
  una sección nueva (§5.1) en la conclusion.
\item Se agregó en la Sección~2 de la tesis una tabla de notación con todas las variables (§2.1.8), y se
  unificó la notación de las fórmulas en toda la tesis.
\item En el apartado 4.1 de la tesis («Configuración experimental») se agregan los tiempos de procesamiento del
  simulador con la máquina actual (Apple M1 Pro con 16~GB de RAM).
\item No cambian el título, los objetivos ni su contenido, el alcance ni la estructura de capítulos. Medido palabra
  por palabra, el 22\,\% del texto es nuevo o modificado; la mayor parte está en los capítulos~4 (45\,\%) y 5
  (41\,\%), donde están las cifras.
\end{itemize}\vspace{2pt}}}

% ─────────────────────────────────────────────────────────────
\section{Criterio y motivos de los cambios}

Para facilitar la nueva revisión se partió del texto revisado y solo se cambió lo que cumple al menos uno de
cuatro motivos. En las secciones~4 y 5 cada cambio lleva la etiqueta de su motivo:

\begin{tabular}{@{}l L{14.2cm}@{}}
\tR & La cifra o la conclusión cambió con la nueva campaña experimental.\\
\tM & La descripción del método debía coincidir con lo que hace el código corregido.\\
\tO & Atiende la observación del comité sobre el planteamiento del problema.\\
\tC & La fuente citada, el código o los datos no sostenían la afirmación.\\
\end{tabular}

Las mejoras de redacción y los análisis adicionales que no cumplen ninguno de estos motivos se dejaron fuera.
Los análisis complementarios (control euclidiano, robustez y tiempo de cómputo) se incluyeron solo en la
extensión necesaria para responder las preguntas de investigación: un párrafo en §4.5.4, una viñeta en §4.5.5
y una oración en §4.1.

Para ver cada cambio en su contexto, el archivo \texttt{tesis\_cambios\_vs\_revision\_2026-06-09.pdf} contiene
la tesis actual con el texto nuevo o modificado resaltado y, en notas, el texto anterior.

% ─────────────────────────────────────────────────────────────
\section{Errores del simulador corregidos}

Al preparar el artículo derivado de la tesis se contrastaron los resultados por semilla con el código. Se
encontraron los errores siguientes; todos se corrigieron antes de la nueva campaña. Ahora cada corrida pasa
verificaciones automáticas: ninguna orden se recoge antes de estar lista ni se entrega antes de recogerse,
ninguna se entrega dos veces y todas las rutas provienen de OSRM.

{\small
\begin{longtable}{@{}L{0.35cm} L{5.1cm} L{5.45cm} L{5.4cm}@{}}
\toprule
\# & Error en la versión revisada & Efecto en los resultados revisados & Corrección (sección actual)\\
\midrule
\endhead
1 & El instante de entrega era el fin del recorrido en OSRM: no incluía la espera de la comida en el
  restaurante ni los 4~min de servicio en la recolección y en la entrega.
  & Entre el 31 y el 70\,\% de las entregas FCFS, y cerca del 2\,\% de las de RH, ocurrían antes de que la
  comida estuviera lista. El CtD medio de FCFS en S1 (10.1~min) era menor que el tiempo medio de preparación
  (12~min).
  & Temporización de Reyes et al.\ (2018): espera hasta que la comida esté lista y mitad del tiempo de servicio
  antes y mitad después de cada recolección y entrega (§3.3.4, nueva).\sep
2 & El despacho solo se ejecutaba cuando el reloj marcaba un minuto múltiplo de 5. El reloj arranca 15~min antes
  de la primera orden, así que si esta llegaba después de las 8:00:59 nunca se despachaba.
  & En las semillas 2027, 2033 y 2034 ninguna orden recibió repartidor. La tesis lo atribuía a una interrupción
  de OSRM y trabajaba con 7 semillas.
  & Despacho en cada época (§3.3.1). Las 10 semillas son válidas.\sep
3 & En RH, el paso de remoción y reinserción reconstruía los agrupamientos a partir de una copia desactualizada:
  una orden movida seguía también en su agrupamiento original.
  & Cerca del 4\,\% de las entregas de RH se contaban dos veces. Esto inflaba el \textit{throughput} de RH
  (su +4.9\,\% en S1 era exactamente $1\,089/1\,038$), la distancia y el costo por pedido.
  & Se reescribió el paso: cada orden queda en un solo agrupamiento, y una entrega doble detiene la simulación
  (§3.5.2).\sep
4 & FCFS tomaba las órdenes cuya comida estaría lista en los siguientes 20~min (el horizonte de RH) y las
  recorría agrupadas por restaurante, no por orden de llegada.
  & La línea base no era la descrita en la tesis («asignación reactiva inmediata»); §3.4.1 reconocía un
  «sesgo de orden de iteración».
  & FCFS despacha cada orden en la primera época posterior a su colocación, en orden de colocación, al
  repartidor libre más cercano (§1.6.2, §3.4.1).\sep
5 & La utilización contaba solo el tiempo de conducción.
  & La lectura «menor utilización, mayor eficiencia» no tenía sustento.
  & Utilización = conducción + servicio, como en Reyes et al.\ (§2.4.3, §3.7.3).\sep
6 & La inserción de órdenes en agrupamientos comparaba el costo absoluto de cada agrupamiento candidato y no el
  incremento de costo.
  & Las órdenes se repartían entre agrupamientos en vez de insertarse donde cuestan menos, a diferencia de
  Reyes et al.\ y de lo descrito en §3.5.2.
  & Se usa el incremento de costo (§3.5.2).\sep
7 & El denominador de $Z_t$ solo incluía repartidores libres.
  & $Z_t$ sobrestimado y agrupamientos más grandes que los definidos en §3.5.1.
  & Repartidores en turno que quedarán libres dentro de $\Delta_2$ (§3.5.1).\sep
8 & El peso de asignación sumaba dos veces el tramo repartidor\,\ra\,restaurante, y el compromiso final estimaba
  la llegada al restaurante como la mitad de la ruta completa.
  & Pesos de asignación y decisiones de compromiso final o parcial incorrectos.
  & Ambos usan los tiempos de §3.3.4, con el tramo OSRM real (§3.5.4).\sep
9 & Tras un compromiso parcial, el repartidor quedaba ubicado en el domicilio del último cliente y no en el
  restaurante.
  & Posición incorrecta del repartidor en las épocas siguientes.
  & El repartidor queda en el restaurante (§3.5.4).\\
\bottomrule
\end{longtable}}

Se corrigieron además dos errores que no afectaban las cifras de la tesis: la caché de rutas no distinguía
entre OSRM y distancia euclidiana, y la opción de sensibilidad del tamaño máximo de agrupamiento no surtía efecto.

Durante la revisión se identificó también una convención del ciclo de simulación que no es un error de
implementación. Como la actualización de rutas ocurre después de la asignación, un repartidor espera entre 5 y
10~min su siguiente ruta. Se mantuvo, se documenta en §3.3.1 y su efecto se evaluó repitiendo la campaña con el
orden inverso (§4.5.5): los signos y el punto de cruce se mantienen, y las magnitudes en S2--S4 se reportan como
cota superior.

% ─────────────────────────────────────────────────────────────
\section{Resultados: antes y ahora}

\subsection{Cifras principales (promedio por escenario)}
Cada celda muestra el valor de la versión revisada (7 semillas) \ra\ el de la versión actual (10 semillas).
$\Delta$CtD es el cambio relativo de RH respecto a FCFS; el SLA es el porcentaje de pedidos entregados en
40~min o menos.

{\small
\begin{longtable}{@{}l c c c c c c@{}}
\toprule
Esc. & $\rho$ & CtD FCFS (min) & CtD RH (min) & $\Delta$CtD & SLA FCFS (\%) & SLA RH (\%)\\
\midrule
S1 & 18 & 10.10 \ra\ 28.93  & 20.38 \ra\ 30.31 & $+101.8$ \ra\ $+4.8$\,\%  & 100.0 \ra\ 84.6 & 99.2 \ra\ 83.4\\
S2 & 26 & 16.68 \ra\ 56.07  & 22.74 \ra\ 35.42 & $+36.3$ \ra\ $-36.8$\,\%  & 91.7 \ra\ 38.1  & 97.0 \ra\ 69.6\\
S3 & 35 & 35.79 \ra\ 108.49 & 25.52 \ra\ 45.36 & $-28.7$ \ra\ $-58.2$\,\%  & 76.3 \ra\ 7.3   & 90.2 \ra\ 57.6\\
S4 & 52 & 75.41 \ra\ 224.88 & 32.26 \ra\ 73.63 & $-57.2$ \ra\ $-67.3$\,\%  & 61.2 \ra\ 6.2   & 78.3 \ra\ 42.5\\
\bottomrule
\end{longtable}}

\subsection{Conclusiones que cambiaron}
{\small
\begin{longtable}{@{}L{2.6cm} L{3.9cm} L{4.95cm} L{4.8cm}@{}}
\toprule
Aspecto & Versión revisada & Versión actual & Por qué cambió\\
\midrule
\endhead
Corridas válidas & 7 semillas, 56 corridas & 10 semillas, 80 corridas & Error 2.\sep
Nivel de los tiempos de entrega & CtD medio de 10 a 75~min & CtD medio de 29 a 225~min
  & Error 1: el CtD incluye ahora la espera de la comida y los tiempos de servicio.\sep
Punto de cruce en CtD & Entre S2 y S3 ($\rho = 26$--35) & Entre S1 y S2 ($\rho = 18$--26)
  & Con la temporización correcta cada viaje ocupa más tiempo al repartidor. FCFS, con un viaje por pedido,
  agota la flota antes; RH comparte la espera y la recolección entre los pedidos de un agrupamiento.\sep
S1 & FCFS 10.3~min más rápido; SLA prácticamente igual
  & FCFS 1.4~min más rápido ($p = 0.0137$); diferencia en SLA no significativa ($p = 0.23$) & Errores 1 a 4.\sep
S2 & FCFS más rápido en CtD; RH mejor en SLA & RH mejor en ambos ($-36.8$\,\% en CtD; $+31.5$~pp en SLA)
  & Errores 1 a 4.\sep
S4 & CtD $-57.2$\,\%; P90 $-77.1$\,\% & CtD $-67.3$\,\%; P90 $-42.0$\,\% & Errores 1 a 4.\sep
Pedidos sin entregar con FCFS & 11.5\,\% en S4 & 11.3\,\% en S3 y 39.9\,\% en S4 (RH: 0 y 0.2\,\%)
  & Error 1. Como el CtD y el SLA se miden sobre pedidos entregados, se agrega el SLA sobre todas las órdenes
  (FCFS 6.4 y 3.8\,\%; RH 57.6 y 42.4\,\%).\sep
\textit{Throughput} & RH superior en todos los escenarios ($+3.4$ a $+17.7$\,\%)
  & Igual en S1--S2; $+12.7$ y $+66.0$\,\% en S3--S4
  & Error 3. Con la misma flota y los mismos turnos, solo difiere por los pedidos que FCFS no entrega.\sep
Utilización & RH menor en S2--S4 & RH menor en S1--S3 y similar en S4 & Errores 1 y 5.\sep
Distancia por pedido & RH mayor en S1 ($+5.3$\,\%) y menor en S2--S4 & RH menor en todos ($-13.5$ a $-49.4$\,\%)
  & Errores 3 y 6.\sep
Tamaño de agrupamiento (RH) & 1.77 a 3.07 & 2.54 a 3.30 por pedido; 1.93 a 2.80 por viaje
  & Errores 6 y 7. Además se aclara que la Tabla~4.5 reporta la media ponderada por pedido.\sep
Costo por pedido en S1 & RH menor (\$11.06 vs \$11.40) & RH mayor (\$11.59 vs \$10.80)
  & Errores 1 y 3. RH concentra las entregas en menos repartidores y el resto cobra el pago mínimo.\sep
Prueba de Wilcoxon & $p = 0.0156$ en las 8 pruebas & $p = 0.0020$ en las 6 pruebas de S2--S4; en S1,
  $p = 0.0137$ (CtD) y $p = 0.23$ (SLA)
  & Con $n = 10$ el $p$ mínimo posible es 0.0020 (con $n = 7$ era 0.0156).\sep
Prueba de Friedman & $\chi^2_F = 21.0$ & $\chi^2_F = 30.0$ ($p < 0.001$) & $n = 10$.\sep
Comparación con Reyes et al. & «Las mejoras superan el 10--25\,\% reportado por Reyes et al.» & Se elimina
  & \tC\ Reyes et al.\ no publican una comparación contra FCFS.\sep
Efecto de la red vial & «La topología de La Paz amplifica la ventaja», sin medición
  & $\kappa = 1.38$ y control euclidiano: la ventaja de RH es mayor con la red real en los cuatro escenarios;
  entre el 62 y el 84\,\% de la diferencia se explica por la mayor longitud media de los viajes
  & \tO\ Responde a PI2 y aporta la evidencia del objetivo específico~2.\sep
Robustez & --- & Las magnitudes en S2--S4 se leen como cota superior & Convención del ciclo (Sección~2).\\
\bottomrule
\end{longtable}}

% ─────────────────────────────────────────────────────────────
\section{Cambios por capítulo}
Las secciones se numeran como en la versión actual. Las figuras 4.1, 4.2 y 4.3 se regeneraron con los datos
nuevos.

\subsection{Capítulo 1. Introducción}
{\small
\begin{longtable}{@{}L{1.75cm} L{8.15cm} L{6.6cm}@{}}
\toprule
Sección & Qué cambió & Por qué\\
\midrule
\endhead
1.1.1 & El TSP es NP-duro porque no se conoce un algoritmo polinomial que lo resuelva; la cita a Lenstra y
  Rinnooy Kan se coloca en esa afirmación.
  & \tC\ Antes se explicaba la NP-dureza por el crecimiento factorial del tiempo, que es un error conceptual.\sep
1.1.2 & Las variantes del VRP (recogida y entrega, entregas divididas) se citan con Cáceres-Cruz et al.; se
  quita «tiempos de servicio variables».
  & \tC\ Laporte (2009) no documenta esas variantes.\sep
1.1.3 & Se quitan «el componente más costoso» y «entre el 30 y el 60\,\% del costo total»; se describe lo que
  revisan Jazemi et al.
  & \tC\ Esas cifras no aparecen en el artículo citado.\sep
1.1.5 & Larsen (2000) se cita como tesis doctoral; se quitan las comillas y «taxonomía fundamental».
  & \tC\ La referencia anterior («Jesper Larsen», artículo en \textit{Computers \& Operations Research}) no
  existe, y el texto entre comillas no era una cita textual.\sep
1.1.6 & Psaraftis (1980) resolvió con programación dinámica el caso de un solo vehículo. Cordeau y Laporte (2007)
  revisan la literatura del DARP.
  & \tC\ Antes se decía que Psaraftis asignaba pasajeros a vehículos y que Cordeau y Laporte «extendieron» su
  trabajo con heurísticas; su artículo es una revisión.\sep
1.1.7 & Ichoua et al.\ estudian el despacho de vehículos en tiempo real con clientes ficticios que representan
  llegadas futuras; las reducciones son 3.4--29.0\,\% (clientes no servidos) y 2.0--16.0\,\% (viaje y retrasos).
  & \tC\ Antes: «\textit{dial-a-ride} dinámico», «información de los próximos 20--30 minutos», «hasta 29\,\%» y
  «4--16\,\%», que no coinciden con el artículo.\sep
1.1.9 & «Cinco características» \ra\ «cuatro»; Reyes et al.\ «argumentan» (antes «evidenciaron»).
  & \tC\ La lista tiene cuatro elementos; el artículo argumenta, no demuestra empíricamente.\sep
1.1.11 & El factor de 1.2--1.4 entre distancia por red y en línea recta se atribuye a Cole y King, citado por
  Boyacı et al. Se recoge su conclusión: sin ventanas de tiempo la aproximación euclidiana es aceptable, pero con
  ventanas estrechas una ruta factible puede volverse infactible. De Pereira et al.\ queda solo el costo de
  licencias.
  & \tC\ Antes: «rutas entre 17 y 40\,\% más largas», «20~min que serían 30» y «frecuentemente infactibles»,
  que Boyacı et al.\ no afirman. «Caja negra» y «auditoría» no están en Pereira et al.\sep
1.2 & Reescrita. Antes tenía tres subsecciones (ineficacia de las aproximaciones tradicionales, hostilidad
  topológica de La Paz, brecha en la eficiencia algorítmica). Ahora tiene cuatro: 1.2.1~el problema de decisión
  (coordinar o reaccionar), 1.2.2~la brecha (ciudades medias y redes reales), 1.2.3~el caso de La Paz y
  1.2.4~preguntas de investigación (pregunta central, PI1--PI3), enlazadas con los objetivos.
  & \tO\ «No se identifica qué problema se ataca»: el texto anterior describía el contexto pero no formulaba el
  problema ni una pregunta. \tC\ También contenía afirmaciones que se corrigen en otras filas: el 17--40\,\%
  atribuido a Boyacı et al.\ y el «96\,\% del óptimo» de Reyes et al.\ (el óptimo de comparación es de Yildiz y
  Savelsbergh, y la cifra es un CtD 4\,\% superior).\sep
1.5.1 & Figliozzi (2007) se cita por lo que muestra (los kilómetros dependen de las restricciones de ruteo y de
  la distancia al área de servicio). La frase de aporte se formula en términos de la pregunta central.
  & \tC\ El artículo no trata sobre el error de las aproximaciones lineales. \tO\ Coherencia con §1.2.\sep
1.6.1 & Los dos picos diarios provienen del perfil de recolección de paquetes de LaDe (9:00 y 17:00) y no
  corresponden a los horarios de comida locales.
  & \tC\ Antes: «adaptados a la zona horaria local». LaDe es un conjunto de datos de paquetería.\sep
1.6.2 & FCFS asigna cada orden en el momento en que se coloca. El escenario de control ejecuta ambas políticas con
  distancias en línea recta a la velocidad media de la red OSRM.
  & \tM\ Error 4. El control descrito antes («FCFS sobre OSRM») es la propia línea base y no aísla el efecto del
  mapa; se describe el control que se ejecutó.\sep
1.6.4 & El tiempo de servicio se aplica con la convención de Reyes et al.\ (mitad antes y mitad después) y el
  repartidor espera si la comida no está lista. La transferibilidad de los patrones de consumo se presenta como
  limitación.
  & \tM\ Error 1. \tC\ Muñoz-Villamizar et al.\ se citaban para afirmar que los patrones son «universales y
  transferibles»; el artículo sostiene que el tipo de ciudad influye.\\
\bottomrule
\end{longtable}}

\subsection{Capítulo 2. Marco teórico}
{\small
\begin{longtable}{@{}L{1.75cm} L{8.15cm} L{6.6cm}@{}}
\toprule
Sección & Qué cambió & Por qué\\
\midrule
\endhead
2.1 & Se explica en qué difieren Yildiz y Savelsbergh (información perfecta) y Ulmer et al.\ (\textit{ready time}
  incierto), y que el simulador no adopta esos elementos.
  & \tC\ Antes se decía que el simulador incorporaba cotas duales y \textit{ready time} aleatorio; no calcula
  cotas, y el despachador conoce el \textit{ready time} desde que se coloca la orden.\sep
2.1.2 & La Ec.~2.1 se presenta como síntesis expositiva; Reyes et al.\ ajustan su algoritmo con un objetivo
  jerárquico (primero entregar casi todas las órdenes, luego minimizar el CtD).
  & \tC\ Reyes et al.\ no plantean una suma ponderada como función objetivo.\sep
2.1.7 & Steever et al.\ estudian en simulación, inspirados en un servicio real, pedidos con platillos de varios
  restaurantes (\textit{virtual food court}).
  & \tC\ Antes: «operadores reales» que «confirman empíricamente».\sep
2.2.1 & $f = 5$~min y horizontes de 20~min corresponden a la variante~7 de la malla de parámetros de Reyes et al.
  & \tC\ Antes: «calibración empírica», un proceso que no se realizó ni se documenta.\sep
2.2.2 & $\beta = 0.1$ está dentro del intervalo $[0, 1]$ que exploran Reyes et al.
  & \tC\ Reyes et al.\ no publican un valor recomendado.\sep
2.2.4 & El problema de asignación se resuelve de forma exacta en tiempo polinomial; las implementaciones modernas
  del Algoritmo Húngaro tienen complejidad $O(n^3)$.
  & \tC\ El algoritmo original de Kuhn no es $O(n^3)$.\sep
2.2.5 & La referencia cruzada apunta a §2.1.6 (variantes del entorno).
  & Antes decía «Sección 2.1.5», que es otra sección.\sep
2.3 & Mismas correcciones de Boyacı et al.\ (factor de Cole y King) y Pereira et al.\ que en §1.1.11.
  & \tC\ Ver §1.1.11.\sep
2.3.2 & Dijkstra con montículos de Fibonacci se presenta como mejora posterior al trabajo original. En lugar de una
  complejidad empírica se da el dato de Geisberger et al.: consultas de menos de 1~ms en la red de Europa
  occidental (18~millones de nodos).
  & \tC\ Los montículos de Fibonacci son posteriores a Dijkstra (1959); la complejidad
  $O(\sqrt{|V|}\log|V|)$ no aparece en Geisberger et al.\sep
2.4.2 & LaDe proviene de la operación de paquetería de Cainiao (jornada de 8:00 a 19:00, picos de recolección a las
  9:00 y 17:00). Los centros del perfil se «fijan», no se «calibran», y no corresponden a los horarios de comida en
  México (alrededor de 14:00 y 20:00).
  & \tC\ Antes se presentaba como demanda de comida («almuerzo y cena») calibrada a horarios locales.\sep
2.4.3 & Hildebrandt y Ulmer, citando estudios previos, señalan que la desviación respecto al tiempo esperado predice
  mejor la satisfacción. La utilización incluye el tiempo de servicio.
  & \tC\ Antes: «demuestran». \tM\ Error 5.\\
\bottomrule
\end{longtable}}

\subsection{Capítulo 3. Desarrollo}
{\small
\begin{longtable}{@{}L{1.75cm} L{8.15cm} L{6.6cm}@{}}
\toprule
Sección & Qué cambió & Por qué\\
\midrule
\endhead
Tabla 3.2 & Tiempo de servicio $s$ \ra\ $s_r$, $s_o$ (por recolección y por entrega).
  & \tM\ Error 1: el código aplica cada uno por separado.\sep
3.2.1 & La intensidad $\lambda(t)$ toma la forma diaria de LaDe; sus centros reproducen la ubicación de los picos de
  paquetería y no los horarios de comida.
  & \tC\ Ver §2.4.2.\sep
3.2.2 & $\mu_{prep} = 12$ y $\sigma_{prep} = 3$ son un supuesto del generador, no calibrado.
  & \tC\ Antes: «seleccionados a partir de rangos observados», sin verificación.\sep
3.2.3 & El rango de 2 a 5~km entre restaurante y cliente es un supuesto del generador.
  & \tC\ Antes: «replicando la distribución de LaDe»; LaDe no reporta distancias restaurante--cliente.\sep
3.3.1 & El despacho ocurre en cada época, sin importar el minuto de inicio. Párrafo nuevo sobre la espera de 5 a
  10~min entre rutas por el orden de las fases.
  & \tM\ Error 2. La convención del ciclo se documenta y su efecto se evalúa en §4.5.5.\sep
3.3.2--3.3.3 & El repartidor acumula tiempo de servicio. En RH la orden cambia de estado al recibir el
  compromiso final.
  & \tM\ Errores 3 y 5.\sep
3.3.4 & Subsección nueva, «Temporización de rutas», con las ecuaciones de recolección y entrega.
  & \tM\ Error 1: describe la corrección central.\sep
3.4.1 & FCFS toma todas las órdenes reveladas sin repartidor, las ordena por instante de colocación y espera en el
  restaurante si la comida no está lista. Se elimina el «sesgo de orden de iteración».
  & \tM\ Error 4.\sep
3.4.2 & «Cota inferior práctica» \ra\ «línea base reactiva». La comparación evalúa las políticas completas y no
  atribuye la mejora a un mecanismo en particular.
  & \tC\ RH difiere de FCFS en varios mecanismos a la vez; el diseño no permite atribuir la mejora a uno.\sep
3.5.1 & El denominador de $Z_t$ es el conjunto de repartidores en turno, libres u ocupados.
  & \tM\ Error 7.\sep
3.5.2 & Se minimiza el incremento de costo; la reinserción mantiene cada orden en un solo agrupamiento.
  & \tM\ Errores 3 y 6.\sep
3.5.3 & SciPy 1.13 resuelve la asignación con una variante del algoritmo de Jonker--Volgenant; se aclara que en el
  documento se le llama Algoritmo Húngaro.
  & \tC\ Antes se afirmaba que SciPy implementa el Algoritmo Húngaro.\sep
3.5.4 & El compromiso final usa la duración OSRM real del tramo repartidor\,\ra\,restaurante y los tiempos de
  §3.3.4; tras un compromiso parcial el repartidor queda en el restaurante.
  & \tM\ Errores 8 y 9.\sep
3.6.2 & Rangos reales de flota por escenario (55--59, 38--41, 29--30 y 19--21 repartidores).
  & \tM\ El número de órdenes varía entre semillas (988 a 1\,047), y con él la flota.\sep
3.6.5 & Subsección nueva, «Escenario de control euclidiano»: distancia en línea recta a 651~m/min y una segunda
  variante a 472~m/min.
  & \tO\ Necesaria para responder PI2 (§4.5.4).\sep
3.7.1 & Se agrega el SLA sobre todas las órdenes generadas; las diferencias de porcentajes se reportan en puntos
  porcentuales.
  & \tR\ El SLA sobre entregados favorece a la política que deja más pedidos sin entregar.\sep
3.7.3, Tabla 3.4 & La utilización incluye el servicio. Los pedidos por repartidor por hora solo difieren por la
  cobertura.
  & \tM\ Error 5. \tC\ Ambas políticas usan la misma flota y los mismos turnos.\\
\bottomrule
\end{longtable}}

\subsection{Capítulo 4. Resultados}
{\small
\begin{longtable}{@{}L{1.75cm} L{8.15cm} L{6.6cm}@{}}
\toprule
Sección & Qué cambió & Por qué\\
\midrule
\endhead
Intro.\ & Diez semillas válidas; se elimina la explicación de la falla de OSRM; se indican las corridas de origen.
  & \tR\ Error 2.\sep
4.1, Tabla 4.1 & Cada semilla genera su propia instancia (988 a 1\,047 órdenes); la tabla muestra rangos de flota;
  80 corridas. Las verificaciones de integridad sustituyen a «cada \texttt{order\_id} aparece una vez». Se indica
  que se usó el binario oficial \texttt{osrm-routed} v5.27.1. Oración nueva sobre el tiempo de cómputo y la
  memoria (FCFS 1.8--2.7~s y RH 57--80~s por día simulado; 194--221~MB).
  & \tR\ \tM\ La verificación anterior no detectaba entregas dobles (error 3). \tO\ El tiempo de cómputo responde
  a PI3.\sep
4.2.1, Tabla 4.2, Fig.~4.1 & Cifras nuevas; columna de SLA en $\Delta$pp. S1: FCFS ligeramente mejor. S2: la
  relación se invierte. Cobertura en S3--S4 y SLA sobre todas las órdenes.
  & \tR\ Ver Sección~3.\sep
4.2.2, Tabla 4.3 & RH tiene menor \textit{Ready-to-Pickup} también en S1. El exceso sobre el SLA crece
  «rápidamente» (antes «exponencialmente»); las reducciones son cercanas al 80\,\% (antes, más del 90\,\%).
  & \tR\ \tC\ No se ajustó ningún modelo exponencial.\sep
4.3.1, Tabla 4.4 & \textit{Throughput} igual en S1--S2; la diferencia en S3--S4 refleja cobertura. Utilización
  menor bajo RH en S1--S3 y similar en S4.
  & \tR\ Errores 3 y 5.\sep
4.3.2 & La distancia por pedido bajo RH es menor en todos los escenarios.
  & \tR\sep
4.4, Tabla 4.5, Fig.~4.2 & Tamaños de agrupamiento nuevos; el pie aclara que \textit{Avg Bundle} es la media
  ponderada por pedido. En S1 el 25.4\,\% de los pedidos viaja solo (antes 54\,\%); en S4 el 80.9\,\% viaja en
  agrupamientos de 3 o 4.
  & \tR\sep
4.4.2 & En S1 el costo por pedido es mayor bajo RH, con su explicación.
  & \tR\sep
4.5.1 & Los cálculos residen en \texttt{20261023\_233115\_osrm}.
  & \tR\ La corrida anterior (\texttt{20260413\_182745}) queda invalidada.\sep
4.5.2, Tabla 4.6 & Wilcoxon con $n = 10$; S1 con $p = 0.0137$ (CtD) y $p = 0.23$ (SLA); Friedman
  $\chi^2_F = 30.0$.
  & \tR\sep
4.5.3 & Baja saturación = solo S1; RH supera a FCFS en S2--S4; el \textit{rank-biserial} pasa de $-0.85$ a $+1.00$.
  & \tR\sep
4.5.4 & Se elimina la comparación con el 10--25\,\% de Reyes et al. Párrafo nuevo sobre la red vial: $\kappa = 1.38$
  (3.30~km por la red frente a 2.39~km en línea recta, sobre 10\,184 trayectos), resultados del control euclidiano
  y fracción explicada por la escala (62--84\,\%).
  & \tC\ Reyes et al.\ no comparan contra FCFS. \tO\ Responde a PI2; sustituye la afirmación sin medición sobre
  «menor conectividad y mayor variabilidad».\sep
4.5.5 & Se quita la viñeta de las tres semillas fallidas. Se agregan: la convención del ciclo y su prueba de
  robustez (reducciones de CtD de 14.6, 41.4 y 63.0\,\% en S2--S4, que se leen como cota superior), el origen de
  paquetería de los picos de LaDe y $s_r = s_o$.
  & \tR\ Error 2 y convención del ciclo. \tC\ LaDe.\sep
Fig.~4.3 & El cruce de las curvas de CtD ocurre entre $\rho = 18$ y 26 (antes, entre 26 y 35).
  & \tR\\
\bottomrule
\end{longtable}}

\subsection{Capítulo 5. Conclusiones}
{\small
\begin{longtable}{@{}L{1.75cm} L{8.15cm} L{6.6cm}@{}}
\toprule
Sección & Qué cambió & Por qué\\
\midrule
\endhead
Intro.\ & El capítulo responde primero las preguntas de §1.2.4. & \tO\sep
5.1 & Sección nueva, «Respuesta a las preguntas de investigación»: pregunta central, PI1, PI2 y PI3, cada una con
  la sección del Cap.~4 que la sostiene.
  & \tO\sep
5.2.1 & 80 corridas (antes 56). & \tR\sep
5.2.2 & OE1: puntos sobre el «área urbana» de La Paz (antes «estructura demográfica real»), 988 a 1\,047 órdenes,
  todas las semillas válidas. OE2: se agregan $\kappa$ y los dos controles euclidianos. OE3: temporización de
  Reyes et al.\ y verificaciones de integridad. OE4: $n = 10$ y valores $p$ nuevos.
  & \tR\ \tM\ El generador no usa datos demográficos; muestrea puntos en el área urbana.\sep
5.3 & Hipótesis: en S1 no se verifica (CtD menor bajo FCFS, SLA sin diferencia significativa, \textit{throughput}
  igual); se verifica en S2--S3; máxima en S4. Nota de cota superior.
  & \tR\ Antes el régimen de baja saturación abarcaba S1--S2 y el \textit{throughput} de RH era mayor por las
  entregas dobles (error 3).\sep
5.4 & Punto de cruce en $\rho \approx 18$--26 (antes 26--35 y «señal temprana en S2»). En S4, $-67.3$\,\% de CtD y
  $-42.0$\,\% de P90. Se elimina «superando el 10--25\,\% de Reyes et al.». La topología se describe con
  $\kappa = 1.38$.
  & \tR\ \tC\ Ver §4.5.4.\sep
5.5 & Se elimina «Robustez de la integración con OSRM». «Pesos ($\beta$, $\theta$)» en lugar de
  «$w_1, \ldots, w_4$». El punto de cruce fino se busca entre $\rho = 18$ y 26.
  & \tR\ La falla no era de OSRM (error 2). \tC\ La función de costo implementada tiene los pesos $\beta$ y
  $\theta$; $w_1, \ldots, w_4$ no existen.\\
\bottomrule
\end{longtable}}

% ─────────────────────────────────────────────────────────────
\section{Bibliografía}
{\small
\begin{longtable}{@{}L{3.6cm} L{6.6cm} L{6.3cm}@{}}
\toprule
Referencia & Qué estaba mal & Corrección\\
\midrule
\endhead
Larsen (2000) & Entrada inexistente: «Jesper Larsen», artículo de \textit{Computers \& Operations Research} con DOI
  inválido & Tesis doctoral de Allan Larsen, Technical University of Denmark (IMM-PHD-2000-73)\sep
Wu et al.\ (LaDe) & Autores, congreso, año y DOI incorrectos; el DOI apuntaba a otro artículo
  & KDD '24, pp.~5991--6002, DOI 10.1145/3637528.3671548\sep
Lenstra y Rinnooy Kan (1981) & El DOI correspondía a otro artículo & DOI 10.1002/net.3230110211\sep
Reyes et al.\ (2018) & Registrado como artículo de la «revista» \textit{Optimization Online}, que es un repositorio
  de preprints & Preprint\sep
Hildebrandt y Ulmer & Año 2021 y sin volumen & 2022, vol.~56, núm.~4, pp.~1058--1084\sep
Jazemi et al.\ (2023) & Número «2413015» & Vol.~13, núm.~24, art.~13015, con DOI\sep
Lin et al.\ (2018) & Coautor «Zhaodong Xu» y DOI incorrecto & «Zhe Xu» y DOI correcto\sep
Psaraftis (1980) & Título inexacto & Título original\sep
Boeing & Volumen, páginas y año & Corregidos\sep
Dantzig y Ramser (1959), Muñoz-Villamizar et al., Ichoua et al. & Sin DOI, número de artículo que no se imprimía,
  guiones tipográficos & Corregidos\\
\bottomrule
\end{longtable}}

% ─────────────────────────────────────────────────────────────
\section{Numeración que cambió}
\begin{itemize}\setlength{\itemsep}{1pt}
\item §1.2: tres subsecciones (1.2.1--1.2.3) \ra\ cuatro (1.2.1--1.2.4), con títulos nuevos.
\item §3.3.4 «Temporización de rutas» es nueva; «Integración con OSRM» pasa de 3.3.4 a 3.3.5. Las ecuaciones del
  Cap.~3 posteriores a ella cambian de número.
\item §3.6.5 «Escenario de control euclidiano» es nueva; «Replicación y análisis de sensibilidad» pasa a 3.6.6 y
  «Almacenamiento de resultados» a 3.6.7.
\item §5.1 «Respuesta a las preguntas de investigación» es nueva; las secciones 5.1--5.4 anteriores pasan a
  5.2--5.5.
\item §2.1.8 «Resumen de notación» es nueva y contiene la Tabla~2.1 (notación); las demás tablas y las figuras
  conservan su número (11 tablas y 5 figuras). La tesis pasa de 85 a 94 páginas.
\end{itemize}

% ─────────────────────────────────────────────────────────────
\section{Lo que se mantuvo sin cambios}
\begin{itemize}\setlength{\itemsep}{1pt}
\item Título, objetivo general y objetivos específicos: se conserva el texto registrado en el protocolo (en el
  objetivo específico~1 solo se agregó la cita de LaDe). La hipótesis conserva lo que afirma; solo se precisó su
  redacción técnica (véase «En resumen»). Su grado de cumplimiento con los resultados nuevos se discute en §5.2 y §5.3.
\item Alcance y delimitación; justificación (§1.5.2--1.5.4 sin cambios y §1.5.1 solo en la cita de Figliozzi y la
  frase de aporte).
\item Estructura de capítulos, diseño experimental (cuatro escenarios, diez semillas, pruebas de Wilcoxon y
  Friedman) y conjunto de métricas.
\item Redacción de las contribuciones y del trabajo futuro, salvo los puntos señalados en la Sección~4.
\end{itemize}

\end{document}
